
Code generation models fail on substantial portions of standard benchmarks. When evaluating model performance, most studies report aggregate pass rates (e.g., pass@k metrics) without examining the composition of failures. Different training objectives may create different optimization pressures that manifest as different failure patterns, but this possibility has not been systematically examined.

This study investigates whether two models trained with different objectives---execution-based reinforcement learning (RL) and instruction tuning---produce different distributions of error types among their failed generations. The comparison uses standard error categories: syntax errors (code fails to parse), runtime errors (code parses but crashes during execution), and assertion errors (code executes completely but produces incorrect output).

The study tests four hypotheses:
\begin{enumerate}
\item Error type distributions differ between the two models (existence test)
\item The RL model produces a higher proportion of assertion errors among failures
\item The RL model produces code that executes deeper before failing
\item The effect persists at fine-grained error taxonomy levels
\end{enumerate}

The models compared differ in multiple ways beyond their training objectives, including architecture, parameter count, and pre-training data. These confounds limit causal interpretation. The results describe observed differences between these specific models rather than establishing general properties of alignment methods.
